\section{Conclusion}
\label{sec:conclusion}

We introduced \rnwise{} Output-Oriented Testing, which inverts the combinatorial
testing paradigm: rather than covering interactions among \emph{inputs} and
observing whatever behaviour results, it constructs a $t$-way covering array over
an abstracted \emph{output-behaviour} space and recovers concrete inputs by
gradient-free inverse mapping. We formalised the Output Covering Array
$\OCA(M;s,q,w)$ and the $\OCov_s$ criterion, established coverage bounds including
$M\ge v^{s}$, and gave a four-stage algorithm that is agnostic to the SUT's
internals.

Empirically, across tabular ML, vision, NLP, and quantum systems, the same
algorithm attains $100\%$ feasible output coverage. On a 30-run tabular study it
uniquely reaches perfect output coverage (Mann--Whitney $U$ vs.\ every baseline,
$p\approx6\times10^{-13}$, $\delta=1.0$), matches the strongest behavioural
baseline on fault detection, and does so at roughly $8\%$ of an exhaustive-input
search budget. The contribution is best understood as a new \emph{adequacy
criterion}---output-interaction coverage---together with a constructive method to
satisfy it, unifying behavioural testing of classical and quantum systems.

\paragraph{Future work}
Three directions stand out: (i) learning or optimising the output partition
$\{\pi_j\}$ automatically, so that abstraction is data-driven rather than
hand-specified; (ii) validating that $\OCov_s$ predicts real defect discovery on
production ML pipelines and on quantum hardware, beyond seeded faults; and (iii)
tightening the covering-array construction with constraint solvers to exploit
output-side feasibility structure, further shrinking suites at higher strengths.

\section*{Data Availability}
The complete implementation, all benchmarks and baselines, the experiment drivers,
and scripts that regenerate every table and figure from raw result artifacts are
available in the accompanying repository, with pinned dependencies and a container
image for deterministic reproduction.
